% !TeX program = xelatex
% 在 docs/qa 目录内编译。需要 TeX Live 2025、XeLaTeX、latexmk 及 Fandol 字体。
% pakperk_qa_build=$(mktemp -d /tmp/pakperk-qa.XXXXXX)
% latexmk -xelatex -interaction=nonstopmode -halt-on-error -file-line-error \
%   -outdir="$pakperk_qa_build" pakperk-functional-testing-zh.tex
% 检查临时目录中的 PDF 后，将其复制到 docs/qa，再清理该临时目录。
\DocumentMetadata{lang=zh-CN}
\documentclass[UTF8,a4paper,zihao=-4,fontset=fandol]{ctexart}
\usepackage{pakperk-qa}
\PakPerkDocument{用户基础测试指南}{1.2}

\newcommand{\ui}[1]{「#1」}
\newcommand{\sampletitle}{\textit{Attention Is All You Need}}
\newcommand{\arxivsample}{\texttt{1706.03762}}
\hypersetup{pdfsubject={iOS 与 Android App 人工功能测试}}

\begin{document}
\PakPerkTopic{启动与论文浏览}
\PakPerkGuideTitle{iOS 与 Android 功能测试}

\begin{PakPerkPanel}
  \PakPerkRow{测试入口}{安装本轮送测的 PakPerk App。从源码安装请参照设备开发指南。\\
    \path{docs/mobile-device-development.md}}
  \PakPerkRow{测试材料}{查找和导入可使用 \sampletitle{}（arXiv \arxivsample{}）。深度阅读使用已处理的样本论文。测试账号、公开评论、待复习条目和未读通知由送测方提供。}
  \PakPerkRow{结果记录}{达到预期后勾选通过。出现偏差或无法执行时，记录用例编号、设备型号、系统版本、App 版本和实际情况。}
\end{PakPerkPanel}

\section{启动与论文浏览}

\PakPerkCase{T01}{启动 App}
  {设备主屏幕上的 PakPerk 图标。}
  {\item 联网后，点按图标打开 App。
   \item 出现分类引导时，点按\ui{暂时跳过}。
   \item 等待\ui{阅读}页面加载。}
  {App 进入\ui{阅读}并显示公开论文。游客可以浏览论文。}

\PakPerkCase{T02}{切换论文}
  {底部导航 → \ui{阅读}。}
  {\item 记下当前论文的标题。
   \item 向上滑动，查看下一篇论文的标题和摘要。
   \item 向下滑动，返回刚才的论文。
   \item 依次打开\ui{文库}和\ui{我的}，再回到\ui{阅读}。}
  {切换论文时，标题、作者和摘要随之更新。向下滑动后显示刚才的论文。从其他页面返回后，阅读页面仍可使用。}

\PakPerkCase{T03}{阅读摘要}
  {\ui{阅读} → 任一公开论文的\ui{摘要}。}
  {\item 打开一篇摘要较长的论文。
   \item 查看标题、作者及 arXiv 编号。
   \item 向下滚动，读到摘要末尾。}
  {页面显示所选论文的标题、作者和 arXiv 编号。摘要可滚动至末尾，文字清晰。}

\clearpage
\PakPerkTopic{深入阅读与智能辅助}
\section{深入阅读与智能辅助}

\PakPerkCase{T04}{阅读引言}
  {样本论文 → \ui{引言}。}
  {\item 打开样本论文，从\ui{摘要}切换到\ui{引言}。
   \item 等待正文加载。加载期间留意页面上的进度或状态提示。
   \item 正文出现后上下滚动，查看不同段落。}
  {加载期间有明确提示。正文按段落显示，滚动时排版稳定。}

\PakPerkCase{T05}{向论文助手提问}
  {样本论文 → \ui{引言} → \ui{助手}。}
  {\item 提交一个可从原文核实的问题，例如「这篇论文提出了什么方法？」。
   \item 打开回答中的一条引用，对照引用所指的原文段落。
   \item 回到助手，再提一个超出论文范围的问题，查看回应。}
  {回答围绕当前论文，引用能打开支撑回答的原文。证据不足时，助手会明确提示。}

\PakPerkCase{T06}{查看关联论文}
  {样本论文 → \ui{关联论文}。}
  {\item 找到一张关系卡片，阅读卡片上的论文标题和关系说明。
   \item 点按卡片，查看打开的论文。
   \item 核对页面标题是否与卡片一致。}
  {卡片清楚说明两篇论文的关系。点按后打开卡片上的论文。}

\clearpage
\PakPerkTopic{来源、离线与设置}
\section{来源、离线与设置}

\PakPerkCase{T07}{打开 arXiv 原文}
  {样本论文工具栏 → arXiv。}
  {\item 在样本论文页点按 arXiv，查看打开的原始页面。
   \item 核对原始页面的论文标题或编号，然后返回 App。
   \item 切换到\ui{引言}，点按\ui{原文}打开 PDF。核对论文标题和版本，再返回 App。}
  {原始页面与 App 中的论文一致。PDF 对应当前论文的版本。从外部页面返回后，App 仍显示原先打开的论文。}

\PakPerkCase{T08}{离线阅读与联网恢复}
  {\ui{阅读}及设备的网络设置。}
  {\item 联网打开一篇公开论文，记下标题。
   \item 关闭 Wi-Fi 和蜂窝数据，再次打开 PakPerk。
   \item 找到刚才的论文，查看标题和摘要。
   \item 恢复网络，刷新\ui{阅读}页面。}
  {断网后仍可查看已缓存的标题和摘要。恢复网络后，阅读页面可以正常刷新。}

\PakPerkCase{T09}{更改语言和外观}
  {\ui{我的} → \ui{设置}。}
  {\item 记下当前的语言和外观设置。
   \item 将\ui{外观}切换为另一种主题，再切换\ui{语言}。
   \item 关闭 App，再次打开并查看界面。
   \item 恢复原来的设置，返回\ui{阅读}。}
  {界面主题和文案随设置更新。再次打开 App 后，仍使用上次选择的语言和外观。}

\clearpage
\PakPerkTopic{账号}
\section{账号}

\PakPerkCase{ACC-01}{登录与退出}
  {\ui{我的} → \ui{登录 / 注册}。}
  {\item 发起登录，在系统浏览器中使用测试账号完成身份验证。
   \item 返回 App，在\ui{我的}核对当前账号，再打开\ui{文库}。
   \item 返回\ui{我的}，点按\ui{退出登录}。
   \item 打开\ui{阅读}并浏览公开论文。}
  {身份验证后可以返回 App，账号与登录时使用的一致，文库可以打开。退出后，\ui{我的}显示游客入口，仍可浏览公开论文。}

\PakPerkCase{ACC-02}{完善资料和修改显示名称}
  {\ui{我的} → \ui{完善资料}。}
  {\item 使用待完善资料的测试账号登录。
   \item 填写未被使用的测试用户名，阅读并同意界面所列条款，完成设置。
   \item 在\ui{我的}点按\ui{修改显示名称}，填写新的名称并保存。
   \item 离开\ui{我的}页面，再次进入并查看资料。}
  {账号资料保存成功。公开用户名与填写的值一致，重新打开页面后可看到更新后的显示名称。}

\PakPerkCase{ACC-03}{申请删除账号}
  {\ui{我的} → \ui{删除账号}。}
  {\item 登录送测方指定的注销测试账号，核对当前账号。
   \item 阅读页面上的删除说明，发起删除申请。
   \item 页面要求重新验证身份时，在系统浏览器完成验证。返回 App 后提交申请。
   \item 查看申请状态。分别尝试打开\ui{文库}和\ui{阅读}。}
  {App 显示申请已受理或正在处理。原会话无法继续访问账号文库，公开论文仍可浏览。}

\clearpage
\PakPerkTopic{搜索与推荐}
\section{搜索与推荐}

\PakPerkCase{DIS-01}{查找并打开论文}
  {\ui{阅读} → \ui{搜索论文} → \ui{查找}。}
  {\item 输入 \sampletitle{} 并搜索。
   \item 将搜索词改为 arXiv 编号 \arxivsample{}，再次搜索。
   \item 在对应结果上点按\ui{打开论文}。}
  {标题和编号均能找到目标论文。打开后显示对应摘要，文库没有新增条目。}

\PakPerkCase{DIS-02}{按主题探索并保存查询}
  {\ui{阅读} → \ui{搜索论文} → \ui{探索}。}
  {\item 输入测试主题，例如 Transformer。选择一个分类后提交搜索。
   \item 查看结果和范围说明，点按\ui{保存查询}。
   \item 退出搜索页。再次进入后，在已保存查询中找到该项。
   \item 移除该查询，查看列表。}
  {搜索结果显示所用主题和分类的覆盖范围。保存的查询可以重新查看，移除后从列表消失。}

\PakPerkCase{DIS-03}{查看推荐原因并反馈}
  {\ui{阅读} → \ui{发现}。}
  {\item 在\ui{文库}确认待阅读队列为空，返回\ui{阅读}并刷新。
   \item 在\ui{发现}中选择\ui{最新}，找到带\ui{原因}入口的推荐卡片。
   \item 打开\ui{原因}并查看说明。
   \item 选择\ui{有帮助}。提交后回到\ui{文库}查看队列。}
  {队列为空时可以进入发现。推荐原因对应所选论文，反馈可以提交，待阅读队列仍为空。}

\clearpage
\PakPerkTopic{文库与待阅读}
\section{文库与待阅读}

\PakPerkCase{LIB-01}{按标题和链接添加同一篇论文}
  {\ui{文库} → \ui{添加论文}。}
  {\item 打开\ui{收件箱}，确认待添加的测试论文尚未出现在列表中。
   \item 按标题搜索 \sampletitle{}，选择对应候选并加入待阅读。
   \item 再次打开\ui{添加论文}，输入 \url{https://arxiv.org/abs/1706.03762} 并确认。
   \item 刷新\ui{收件箱}，查看这篇论文的条目数。}
  {按标题查找时先显示候选。确认添加后，收件箱出现目标论文。再次通过链接添加，收件箱中只有一条对应记录。}

\PakPerkCase{LIB-02}{保存、同步及移除论文}
  {论文页工具栏的待阅读按钮及\ui{文库}。}
  {\item 打开一篇尚未收藏的论文，加入待阅读。
   \item 打开\ui{文库}并下拉刷新，等待同步完成。
   \item 退出登录，再使用同一账号登录，核对已保存的论文。
   \item 在文库中移除该论文，再次下拉刷新。}
  {保存后可在文库看到论文。完成同步后，重新登录同一账号仍能看到它。移除后，列表中不再显示该论文。}

\PakPerkCase{LIB-03}{验证待阅读优先和归档}
  {\ui{阅读}及\ui{文库}。}
  {\item 使用只有一篇待阅读论文的测试账号打开\ui{阅读}，查看队列。
   \item 打开这篇论文。返回\ui{文库}，核对其当前状态。
   \item 将它移至\ui{已归档}，等待同步完成。
   \item 刷新\ui{阅读}查看\ui{发现}，再回到文库查找已归档论文。}
  {队列有论文时，\ui{阅读}显示待阅读内容。打开论文不会改变文库状态。最后一篇归档并同步后，阅读页恢复发现，论文仍可在归档中找到。}

\clearpage
\PakPerkTopic{阅读方式与原文定位}
\section{阅读方式与原文定位}

\PakPerkCase{RD-02}{切换阅读模式并恢复进度}
  {样本论文 → \ui{引言}。}
  {\item 依次切换\ui{速览}、\ui{阅读}和\ui{精读}，查看模式及论文标题。
   \item 回到\ui{阅读}模式，打开\ui{大纲}并选择一个章节。
   \item 到达章节后，使用阅读进度控件保存当前位置。
   \item 退出论文。重新打开同一篇论文并恢复阅读进度。}
  {三种模式显示同一篇论文。大纲跳到选定章节，恢复阅读时回到对应位置。文库状态保持原样。}

\PakPerkCase{RD-03}{查看术语与图表原文}
  {样本论文 → \ui{引言} → \ui{精读}。}
  {\item 打开语义提示，点按一个标出的术语，查看释义及来源。
   \item 从\ui{图、表与公式}中打开一个条目，查看编号和说明。
   \item 点按条目的原文入口，核对打开的论文和版本。页面给出页码时，核对 PDF 中对应的页。}
  {术语释义标明来源。图、表或公式的编号和说明与源论文相符。原文入口打开对应版本，可定位的条目跳到正确页。}

\PakPerkCase{RD-04}{核对论文概览的来源}
  {样本论文\ui{摘要} → \ui{论文概览} → \ui{查看完整论文概览}。}
  {\item 打开完整概览，核对论文标题。
   \item 选择一项附有来源的概览内容，打开对应证据。
   \item 对照证据中的原文，核对概览说法。}
  {概览属于当前论文。证据能定位到相关原文，概览中的说法与所引内容一致。}

\clearpage
\PakPerkTopic{私密研究}
\section{私密研究}

\PakPerkCase{RES-01}{保存并重新打开本机笔记}
  {论文页侧边的\ui{本篇论文的笔记}图标。}
  {\item 打开笔记，输入一句测试文字。
   \item 看到\ui{已保存在本机}后关闭笔记页。
   \item 在同一设备上重新打开这篇论文的笔记。}
  {保存的笔记仍可查看，文字与输入一致，并显示在这篇论文下。}

\PakPerkCase{RES-02}{从原文添加批注并定位}
  {样本论文 → \ui{引言}。}
  {\item 长按原文并选中一句话。
   \item 点按\ui{添加笔记}，输入测试文字并保存。
   \item 打开\ui{研究} → \ui{笔记}，找到刚保存的条目。
   \item 点按该条目的原文入口。}
  {笔记列表显示刚保存的批注。打开原文时，阅读器回到所选文字的位置。文库状态保持原样。}

\PakPerkCase{RES-03}{推迟待复习条目}
  {\ui{我的} → \ui{研究记忆}。}
  {\item 找到测试账号中的待复习条目，点按\ui{打开来源论文}。
   \item 核对论文标题，然后返回\ui{研究记忆}。
   \item 点按该条目的\ui{推迟}，选择\ui{明天}。
   \item 刷新待复习列表。}
  {来源入口打开对应论文。推迟后，该条目从当前待复习列表移出，来源论文的文库状态保持原样。}

\clearpage
\PakPerkTopic{阅读更新}
\section{阅读更新}

\PakPerkCase{ENG-01}{生成并继续阅读简报}
  {\ui{我的} → \ui{设置} → \ui{阅读更新} → \ui{今日阅读简报}。}
  {\item 在文库中确认有待阅读论文，打开\ui{今日阅读简报}并点按\ui{生成今日简报}。
   \item 查看第一篇论文和简报进度，点按\ui{继续}。
   \item 离开页面，再次进入\ui{今日阅读简报}，查看当前位置。
   \item 到\ui{文库}核对简报中的论文。}
  {简报使用待阅读队列中的论文。重新打开后，简报停在上次读到的位置，文库中的论文和状态保持原样。}

\PakPerkCase{ENG-02}{添加订阅并标记通知已读}
  {\ui{我的} → \ui{设置} → \ui{阅读更新}。}
  {\item 点按\ui{添加订阅}，创建一个分类订阅。
   \item 在\ui{应用内通知}中找到测试账号的未读通知。
   \item 点按\ui{标记已读}。
   \item 离开页面，再次进入，核对订阅列表和通知状态。}
  {订阅列表显示刚创建的分类订阅。标记过的通知保持已读。}

\clearpage
\PakPerkTopic{讨论与问卷}
\section{讨论与问卷}

\PakPerkCase{COM-01}{浏览讨论并管理自己的评论}
  {论文页工具栏的评论图标。个人评论位于\ui{我的} → \ui{我的评论}。}
  {\item 以游客身份打开样本论文的讨论，查看公开评论。
   \item 登录测试账号，回到同一篇论文，发送一条测试评论。
   \item 从\ui{我的评论}找到这条评论，回到讨论页并修改文字。
   \item 使用另一个会话查看这篇论文的公开讨论，核对评论的可见性。
   \item 切回作者会话，删除测试评论，再刷新讨论页和\ui{我的评论}。}
  {游客可以查看已发布的评论。作者能找到、编辑并删除自己的评论。已发布的测试评论可在公开讨论中看到，审核中的测试评论仅作者可见。删除后，列表中不再显示该评论。}

\PakPerkCase{SUR-01}{保存并提交体验问卷}
  {\ui{我的} → \ui{聊聊使用感受}。}
  {\item 以游客身份打开问卷，先回答一到两题。
   \item 返回\ui{我的}。重新打开问卷，查看已选答案。
   \item 完成剩余题目，点按\ui{提交}。}
  {重新打开问卷后保留已选答案。游客可以提交问卷，提交成功后页面显示感谢提示。}

\end{document}
